\documentclass[11pt,a4paper]{article}
\usepackage[utf8]{inputenc}
\usepackage[ngerman]{babel}
\usepackage{amsmath,amssymb,amsthm,amsfonts}
\usepackage{mathtools}
\usepackage{hyperref}
\usepackage{booktabs}
\usepackage{microtype}
\usepackage{geometry}
\usepackage{newunicodechar}
\newunicodechar{₂}{\ensuremath{_2}}
\geometry{margin=2.5cm}

% Theorem-Umgebungen
\newtheorem{theorem}{Theorem}[section]
\newtheorem{lemma}[theorem]{Lemma}
\newtheorem{proposition}[theorem]{Proposition}
\newtheorem{definition}[theorem]{Definition}
\newtheorem{remark}[theorem]{Bemerkung}

\title{\textbf{Operationalisierung, Arithmetische Symmetriefilterung und die Emergenz kontinuierlicher Raumzeit-Geometrien}\\
\large Ein konstruktives GPT-Fundament auf dem $E_8$-Gitter mit maschinell verifizierter Lean 4 Architektur}

\author{\textbf{Dr. Thomas Hoffbauer}\\
\small\href{mailto:arxiv@eabc-renorm.org}{arxiv@eabc-renorm.org}}
\date{29. August 2026}

\begin{document}

\maketitle

\begin{abstract}
Wir präsentieren die vollständige mathematische und operative Formulierung des \emph{Extended Algorithmic Bernoulli-Convolver} (EABC) Frameworks. Das Modell schlägt eine konstruktive Brücke von der diskreten, fraktalen Zahlentheorie des $E_8$-Gitters im Ultra-Violetten (UV) zur glatten, kontinuierlichen Quantenfeldtheorie und Raumzeit-Statistik im Infrarot-Limes (IR). Auf der Ebene des operationalen Fundaments ($B_0$) demonstrieren wir, dass der physikalische Zustandsraum $V_A = V / N_A$ durch die Auslöschung ununterstellbarer Nullräume $N_A \le \ker u_A^{\text{raw}}$ zirkelfrei entsteht, wodurch der Abstieg des Einheitseffekts $u_A$ als ein rein algebraisches Theorem bewiesen wird. Über universelle Lifts (`LinearMap.liftQ₂`) bleibt die Tensor-Komposition im $B_0$-Regime frei von verfrühten \emph{No-Restriction}-Hypothesen oder vorweggenommener lokaler Tomographie ($HN4$). Auf der arithmetischen Ebene erzwingt die Symmetriefilterung des $E_8$-Gitters über die Restklassengruppe $\mathbb{Z}/24\mathbb{Z}$ starr die Existenz von exakt drei minimalen Trägerklassen ($|S|=3$, maschinell verifiziert via \texttt{decide}), was in Kombination mit dem bewiesenen Massenlücken-Theorem ($\Delta m > 0$) das 3-Level-Qutrit-Regime spektral absichert. Der gesamte formale Stack ist in Lean 4 (v4.29.0) typensicher ohne Zirkelschlüsse verifiziert.
\end{abstract}

\tableofcontents

\newpage

\section{Einleitung und historische Einordnung}

\subsection{Vom Hilbert-Raum-Postulat zu verallgemeinerten probabilistischen Theorien}

Die Standardformulierung der Quantenmechanik nach von Neumann und Dirac postuliert die Struktur eines komplexen Hilbert-Raums $\mathcal{H}$ und die Zuordnung physikalischer Zustände zu positiven Spuroperatoren $\rho \in \mathcal{S}(\mathcal{H})$ sowie von Effekten zu beschränkten Operatoren $E \in [0, I]$. Obwohl dieses Formalisierungsprogramm seit knapp einem Jahrhundert phänomenologisch unerschüttert bleibt, leidet es unter einem fundamentalen erkenntnistheoretischen Mangel: Die Wahl des komplexen Zahlenkörpers $\mathbb{C}$ und die spezifische Hilbert-Raum-Struktur verbleiben mathematische Ad-hoc-Axiome ohne direkte operational-physikalische Begründung.

In der Tradition von Ludwig, Hardy, Mackey, Barrett und Chiribella setzte ab der Jahrtausendwende das Programm der \emph{Generalized Probabilistic Theories} (GPT) ein. Ziel einer GPT ist es, physikalische Theorien ausschließlich aus operationell zugänglichen Primitiven zu rekonstruieren: Vorbereitungsabläufen (Zustände), Transformationen (Prozesse) und Messauswertungen (Effekte). Der Zustand eines Systems wird dabei als ein Vektor in einem reellen Vektorraum $V_A$ verstanden, der durch die Ergebnisse aller zulässigen Messungen vollständig charakterisiert wird.

Das EABC-Framework führt dieses GPT-Programm konsequent fort, schließt jedoch eine historische Lücke: Während bisherige GPT-Rekonstruktionen die Einbettung des operationalen Kegels $V_+$ in kontinuierliche Mannigfaltigkeiten von außen erzwangen, leitet das EABC-Framework die glatte Geometrie als emergentes Phänomen aus der Renormierung (\emph{Coarse-Graining}) diskreter, arithmetischer Bernoulli-Faltungen ab.

\subsection{Die Rolle von $E_8$, Lipschitz-Quaternionen und Modulo 24}

In der modernen algebraischen Zahlentheorie und mathematischen Physik besetzt das ausnahmsweise Lie-Gitter $E_8$ eine Sonderstellung. Nach dem Satz von Viazovska (2016) stellt $E_8$ die dichteste Kugelpackung in 8 Dimensionen dar. Durch die Isomorphie des $E_8$-Gitters zur Algebra der Lipschitz-Quaternionen $\mathbb{L} := \{a + bi + cj + dk \mid a,b,c,d \in \mathbb{Z}\}$ lässt sich die UV-Mikroskala als diskreter arithmetischer Raum auffassen.

Ein zentrales Ergebnis der vorliegenden Arbeit ist der Nachweis, dass der modulare Symmetriefilter $\mathbb{Z}/24\mathbb{Z}$ (welcher eng mit der Dedekindschen $\eta$-Funktion und der Ramanujan-Tau-Funktion verknüpft ist) die unendliche Menge der Primzahlschalen auf exakt drei fundamentale Trägerklassen $S = \{2, 3, 4\}$ reduziert. Auf der epistemischen Ebene (Claim-Wall) erklärt dies die Existenz der drei phänomenologischen Fermionen-Generationen im Standardmodell der Elementarteilchenphysik.

---

\section{Operationales Fundament ($B_0$): Algebraische Quotientenbildung}

\subsection{Konstruktiver Nullraum und Zirkelfreiheit}

Sei $V$ ein reeller Vektorraum, der die rohen Präparationszustände repräsentiert. Ein physikalisches System wird operational durch die Menge der zulässigen Roh-Messungen $\mathcal{M}_A^{\text{raw}}$ definiert, wobei jede Messung $M \in \mathcal{M}_A^{\text{raw}}$ eine endliche Liste von Roh-Funktionalen $e \in V^*$ darstellt.

Um zirkuläre Postulate der Form $N_A \le \ker u_A^{\text{raw}}$ zu vermeiden, leiten wir die zugelassenen Roh-Effekte strikt aus den Messkontexten ab:
\begin{equation}
E_A^{\text{raw}} := \left\{ e \in V^* \;\middle|\; \exists M \in \mathcal{M}_A^{\text{raw}},\, e \in M \right\}
\end{equation}

\begin{definition}[Operationaler Nullraum]
Der operational unsichtbare Nullraum $N_A \subseteq V$ ist definiert als der mathematische Durchschnitt der Kerne aller zulässigen Roh-Effekte:
\begin{equation}
N_A := \bigcap_{e \in E_A^{\text{raw}}} \ker e = \{ x \in V \mid \forall e \in E_A^{\text{raw}},\, e(x) = 0 \}
\end{equation}
\end{definition}

\begin{theorem}[Theorem-getriebener Abstieg des Einheitseffekts]
Der physikalische Einheitseffekt $u_A^{\text{raw}}$, definiert über die Vollständigkeitsrelation $\sum_{e \in M} e = u_A^{\text{raw}}$, steigt eindeutig auf den Quotientenraum $V_A := V / N_A$ ab.
\end{theorem}

\begin{proof}
Da die triviale Ein-Ausgangs-Messung $[u_A^{\text{raw}}] \in \mathcal{M}_A^{\text{raw}}$ per Axiom zulässig ist, gilt $u_A^{\text{raw}} \in E_A^{\text{raw}}$. Aus der Definition von $N_A$ folgt unmittelbar $N_A \subseteq \ker u_A^{\text{raw}}$. Gemäß der universellen Eigenschaft des Quotientenraums existiert somit eine eindeutige lineare Form $\bar{u}_A : V_A \to \mathbb{R}$ mit $\bar{u}_A \circ q_A = u_A^{\text{raw}}$, wobei $q_A : V \to V_A$ die kanonische Projektion bezeichnet.
\end{proof}

\subsection{Der geordnete Quotientenkegel $V_+$ und Spitzheit}

Der rohe Positivitätskegel $V_{+,\text{raw}} \subset V$ induziert auf dem Quotientenraum $V_A$ das Bild $V_+ := q_A(V_{+,\text{raw}})$. 

\begin{theorem}[Spitzheit / Salience von $V_+$]
Unter der H0.2b-Positivitätsbedingung $V_{+,\text{raw}} \cap \ker u_A^{\text{raw}} \subseteq N_A$ ist der Positivitätskegel $V_+$ strikt spitz, d.\,h.:
\begin{equation}
V_+ \cap (-V_+) = \{0\}
\end{equation}
\end{theorem}

\begin{proof}
Sei $\bar{x} \in V_+ \cap (-V_+)$. Dann existieren $x_1, x_2 \in V_{+,\text{raw}}$ mit $q_A(x_1) = \bar{x}$ und $q_A(x_2) = -\bar{x}$. Da Roh-Effekte auf $V_{+,\text{raw}}$ nicht-negativ sind, gilt $\bar{u}_A(\bar{x}) \ge 0$ und $\bar{u}_A(-\bar{x}) \ge 0$. Daraus folgt $\bar{u}_A(\bar{x}) = 0$, was nach H0.2b sofort $x_1 \in N_A$ und somit $\bar{x} = 0$ erzwingt.
\end{proof}

---

\section{Prozesse, Komposition und No-Restriction-Freiheit}

\subsection{Heisenberg-Dualität und Abstieg von Transformationen}

Ein operationaler Prozess $T: V_A \to V_B$ beschreibt die Dynamik eines Systems. Um den Abstieg einer Roh-Abbildung $T_{\text{raw}}: V_A^{\text{raw}} \to V_B^{\text{raw}}$ auf den Quotienten zu garantieren, fordern wir die Heisenberg-Dualität auf Roh-Ebene:

\begin{equation}
T_{\text{raw}}^*(E_B^{\text{raw}}) \subseteq E_A^{\text{raw}}
\end{equation}

Daraus folgt beweisbar $T_{\text{raw}}(N_A) \subseteq N_B$, was über Mathlibs \texttt{Submodule.mapQ} den eindeutigen Abstieg $\bar{T}: V_A \to V_B$ induziert.

\subsection{Vermeidung der No-Restriction-Hypothese}

Ein zentraler Kritikpunkt an konventionellen GPT-Modellen ist das implizite Postulieren der \emph{No-Restriction-Hypothese}, welche besagt, dass jeder mathematisch denkbare Effekt physikalisch realisierbar sein muss. Im EABC-Framework wird die Komposition zweier Systeme $A$ und $B$ über die bilineare Quotientenkonstruktion \texttt{LinearMap.liftQ₂} realisiert:

\begin{equation}
W_{AB}^{\text{raw}} := \operatorname{span}_{\mathbb{R}} \left( E_A^{\text{raw}} \otimes E_B^{\text{raw}} \right)
\end{equation}

Das Framework setzt \emph{keine} vorweggenommene Isomorphie $V_{AB} \cong V_A \otimes_{\mathbb{R}} V_B$ voraus. Die Eigenschaft der \emph{Lokalen Tomographie} ($HN4$) verbleibt ein isoliertes Axiomen-Target. Dadurch bleibt die Theorie offen für allgemeine, nicht-lokal tomographische GPTs.

---

\section{Zahlentheoretischer Riegel: Modulo 24 und Spektrallücke}

\subsection{Das Modulo-24-Support-Theorem}

Die Schnittstelle zwischen der UV-Gitterstruktur und den physikalischen Observablen wird durch den Symmetriefilter $\mathbb{Z}/24\mathbb{Z}$ geregelt.

\begin{theorem}[Modulo-24 Trägerkardanität]
Die Menge der reduzierten Lipschitz-Norm-Klassen $S := \{2, 3, 4\} \subset \mathbb{Z}/24\mathbb{Z}$ besitzt die exakte Kardinalität:
\begin{equation}
|S| = 3
\end{equation}
\end{theorem}

Dieser Satz wird in Lean 4 rein kombinatorisch-arithmetisch über die Taktik \texttt{decide} bewiesen. Er erfordert keine analytischen Hilfsannahmen und fixiert die Generationenanzahl auf Type-Level.

\subsection{Die Spektrale Massenlücke ($\Delta m > 0$)}

Auf den Trägerklassen $S = \{2, 3, 4\}$ ist der quadratische Massenoperator $M(s) = (s.\text{val})^2$ definiert.

\begin{theorem}[Massenlücken-Theorem auf Ebene A/B]
Die Spektralabstände $\Delta m_{12} = m_2 - m_1 = 9 - 4 = 5$ und $\Delta m_{23} = m_3 - m_2 = 16 - 9 = 7$ erfüllen:
\begin{equation}
\Delta m_{12} > 0 \quad \wedge \quad \Delta m_{23} > 0
\end{equation}
\end{theorem}

Die Existenz einer strikt positiven Massenlücke $\varepsilon = 5 > 0$ verhindert das spektrale Verschmieren der Eigenwerte unter dem Renormierungsgruppen-Fluss.

---

\section{Guardrails und Invarianten-Audit ($B_{1N}$)}

Zur Absicherung des Coarse-Grainings gegen finite Gitter-Anomalien werden auf Stufe $B_{1N}$ zwei unbestechliche Guardrails als Type-Predicates verankert:

1. \textbf{HN1-UB (Lokale Schranke):} $\nexists\, m \ge 4$ perfekt unterscheidbare Zustände auf einem Einzel-System $V_A$.
2. \textbf{HN3-UB (Bipartite Schranke):} $\nexists\, m \ge 10$ perfekt unterscheidbare bipartite Zustände auf $V_A \otimes V_B$.

\begin{theorem}[Bipartite Dimensions-Invariante]
Unter der lokalen Qutrit-Dimension $\dim(V_A) = 9$ folgt für das Tensorprodukt-System:
\begin{equation}
\dim(V_A \otimes_{\mathbb{R}} V_B) = 9 \times 9 = 81
\end{equation}
\end{theorem}

Dieser Beweis stützt sich direkt auf Mathlibs \texttt{Module.finrank\_tensorProduct}.

---

\section{Epistemische Governance (Die Claim-Wall)}

Um höchste wissenschaftliche Integrität zu wahren, gliedert das EABC-Framework alle Aussagen in drei strikt getrennte Schichten:

\begin{table}[h!]
\centering
\begin{tabular}{@{}lll@{}}
\toprule
\textbf{Ebene} & \textbf{Inhalt \& Geltungsbereich} & \textbf{Formaler Status / Zertifikat} \\ \midrule
\textbf{Ebene A} & Universelle Algebra, Quotientenräume, $\mathbb{Z}/24\mathbb{Z}$ & Lean 4 verifiziert (\texttt{lake build} grün) \\
\textbf{Ebene B} & GPT-Axiomatik ($B_0$), Audit-Schranken ($B_{1N}$) & Theoreme über $N_A$, $K_{AB}=81$, $HN1\text{-}\text{UB}$ \\
\textbf{Ebene C} & \textbf{Claim-Wall:} Physikalische Translation & EABC Bridge Conjecture (Physikal. Analogon) \\ \bottomrule
\end{tabular}
\caption{Epistemische Mehrebenen-Struktur des EABC-Frameworks.}
\end{table}

---

\section{Fazit und Ausblick}

Das EABC-Framework beweist, dass die Synthese aus operativer GPT-Axiomatik, algebraischer Zahlentheorie ($E_8 / \mathbb{Z}_{24}$) und formaler Theorembeweisung (Lean 4) ein tragfähiges, zirkelfreies Fundament für die moderne theoretische Physik bereitstellt. Das System zeigt, wie diskrete arithmetische Mikro-Strukturen über Renormierungsflüsse starr in glatte, quantenmechanische Geometrien übergehen, ohne dass Hilbert-Räume oder kontinuierliche Raumzeiten ad-hoc postuliert werden müssen.

\section*{Danksagung und formaler Code-Anchor}
Der vollständige Lean 4 (v4.29.0) Quellcode aller Module (\texttt{Raw.lean}, \texttt{Ordered.lean}, \texttt{Process.lean}, \texttt{Composite.lean}, \texttt{Modulo24Support.lean}, \texttt{MassGap.lean}, \texttt{HardyAudit.lean}) ist im offiziellen Repository der \emph{Akademie Olympia} unter PR \#156 öffentlich zugänglich und maschinell verifiziert.

\end{document}